\documentclass[12pt]{article}

\usepackage[utf8]{inputenc}
\usepackage[spanish, es-noshorthands]{babel}

\usepackage{mathtools}
\usepackage{amssymb}
\usepackage{amsthm}

% --- Resultados (cursiva)
\theoremstyle{plain}
\newtheorem{theorem}{Teorema}
\newtheorem{lemma}{Lema}
\newtheorem{proposition}{Proposición}
\newtheorem{corollary}{Corolario}

% --- Definiciones / ejemplos (texto normal)
\theoremstyle{definition}
\newtheorem*{definition}{Definición}
\newtheorem*{example}{Ejemplo}
\newtheorem{algorithmthm}{Algoritmo}

% --- Observaciones (remark)
\theoremstyle{remark}
\newtheorem{remark}{Observación}
\newtheorem{note}{Nota}

\usepackage{geometry}
\geometry{margin=2.5cm}

\usepackage{tikz}
\usetikzlibrary{arrows.meta, positioning, calc, decorations.markings}
\tikzset{
  v/.style={circle, draw, fill=black, inner sep=1.6pt},
  every edge/.style={draw, thick}
}

\usepackage{hyperref}

\title{Teoría de Grafos}
\author{Juan Felipe Rodríguez Córdoba}
\date{}

\begin{document}

\maketitle
\section*{Grafo}
\subsection*{Definiciones}
    \begin{itemize}
        \item Un grafo es una pareja de conjuntos $G = (V, A)$ donde $V = \{a, b, c, \dots\}$ es el conjunto de vértices y $A = \{(a, c), (a, b), \dots\}$ es el conjunto de aristas; parejas no ordenadas de vértices.
        \item Si dos vertices $a$ y $b$ están unidos por multiples aristas, se dice que el grafo es multigrafo.
        \item Si el grafo tiene una arista que une un vértice consigo mismo, esta arista se llama lazo y el grafo se llama pseudografo.
        \item Un grafo es simple si no tiene lazos ni aristas múltiples.
    \end{itemize}

\subsection*{Tipos de grafos}

\begin{itemize}
    \item \textbf{Grafo simple}: no tiene lazos ni aristas múltiples.
    \item \textbf{Multigrafo}: puede tener aristas paralelas (varias aristas entre dos vértices).
    \item \textbf{Pseudografo}: puede tener lazos.
    \item \textbf{Grafo dirigido (digrafo)}: las aristas tienen orientación.
    \item \textbf{Grafo ponderado}: las aristas tienen pesos asociados (distancia, coste, etc.).
\end{itemize}

\begin{figure}[htbp]
\centering
\begin{tikzpicture}[scale=1]
% --- Simple ---
\node at (-3.2,1.4) {\small Simple};
\node[v] (a1) at (-4,0) {};
\node[v] (b1) at (-2.4,0.6) {};
\node[v] (c1) at (-2.4,-0.6) {};
\draw (a1)--(b1)--(c1)--(a1);

% --- Multigrafo ---
\node at (0,1.4) {\small Multigrafo};
\node[v] (a2) at (-0.8,0) {};
\node[v] (b2) at (0.8,0) {};
\draw (a2) .. controls (0,0.8) .. (b2);
\draw (a2) -- (b2);
\draw (a2) .. controls (0,-0.8) .. (b2);

% --- Pseudografo ---
\node at (3.2,1.4) {\small Pseudografo};
\node[v] (a3) at (2.4,0) {};
\node[v] (b3) at (4.0,0) {};
\draw (a3) -- (b3);
\draw (a3) .. controls (2.0,0.6) and (2.0,-0.6) .. (a3); % lazo
\end{tikzpicture}
\caption{Ejemplos: grafo simple, multigrafo (aristas paralelas) y pseudografo (lazo).}
\end{figure}

\subsection*{Representaciones de un grafo}

Un mismo grafo puede representarse de muchas formas distintas sin cambiar su estructura.

Dos representaciones se llaman \textbf{equivalentes} o \textbf{isomorfas} si describen el mismo conjunto de vértices y aristas.

\subsection*{Isomorfismo de grafos}

Dos grafos son isomorfos si existe una biyección entre sus conjuntos de vértices que preserva las aristas.

Es decir, si dos vértices están conectados en un grafo, sus correspondientes también lo están en el otro.

Para verificar si dos grafos son isomorfos se suelen comparar:

\begin{itemize}
    \item Número de vértices.
    \item Número de aristas.
    \item Grados de los vértices.
    \item Existencia de ciclos.
    \item Cantidad de caminos de cierta longitud.
\end{itemize}

\subsection*{Grafo completo}

Un grafo con $n$ vértices se llama \textbf{completo} (denotado $K_n$) si cada pareja de vértices está conectada por una arista.

El número de aristas de $K_n$ es:

\[
\binom{n}{2} = \frac{n(n-1)}{2}
\]

\begin{figure}[htbp]
\centering
\begin{tikzpicture}[scale=1.2]
\node[v] (1) at (90:1.4) {};
\node[v] (2) at (18:1.4) {};
\node[v] (3) at (-54:1.4) {};
\node[v] (4) at (-126:1.4) {};
\node[v] (5) at (162:1.4) {};
\foreach \i in {1,2,3,4,5}{
  \foreach \j in {1,2,3,4,5}{
    \ifnum\i<\j
      \draw (\i)--(\j);
    \fi
  }
}
\node at (0,-2.0) {\small $K_5$};
\end{tikzpicture}
\caption{Grafo completo $K_5$: cada pareja de vértices está conectada.}
\end{figure}

\subsection*{Grafo bipartito}

Un grafo es \textbf{bipartito} si su conjunto de vértices puede dividirse en dos subconjuntos disjuntos $V_1$ y $V_2$ tales que no hay aristas entre vértices del mismo subconjunto.

El grafo bipartito completo se denota $K_{n,m}$ y conecta todos los vértices de $V_1$ con todos los de $V_2$.

Su número de aristas es:

\[
|A| = n \cdot m
\]

\begin{figure}[htbp]
\centering
\begin{tikzpicture}[scale=1]
\node[v,label=left:$u_1$] (u1) at (0,1) {};
\node[v,label=left:$u_2$] (u2) at (0,0) {};
\node[v,label=left:$u_3$] (u3) at (0,-1) {};

\node[v,label=right:$v_1$] (v1) at (3,1) {};
\node[v,label=right:$v_2$] (v2) at (3,0) {};
\node[v,label=right:$v_3$] (v3) at (3,-1) {};

\foreach \u in {u1,u2,u3}{
  \foreach \v in {v1,v2,v3}{
    \draw (\u)--(\v);
  }
}
\node at (1.5,-1.7) {\small $K_{3,3}$};
\end{tikzpicture}
\caption{Bipartito completo $K_{3,3}$: todas las conexiones entre dos partes.}
\end{figure}

\subsection*{Grado de un vértice}

El grado de un vértice es el número de aristas incidentes en él.

En grafos con lazos, cada lazo cuenta como dos unidades de grado.

Para un grafo simple se cumple:

\[
2|A| = \sum_{v \in V} \mathrm{gr}(v)
\]

Es decir, la suma de los grados de todos los vértices es el doble del número de aristas.

\subsection*{Caminos y conectividad}

\begin{itemize}
    \item \textbf{Camino}: sucesión de vértices consecutivos conectados por aristas.
    \item \textbf{Camino simple}: no repite vértices.
    \item \textbf{Recorrido}: no repite aristas.
    \item \textbf{Circuito}: recorrido cerrado.
    \item \textbf{Ciclo}: circuito que no repite vértices.
    \item \textbf{Grafo conexo}: existe un camino entre cualquier par de vértices.
\end{itemize}

\subsection*{Árboles}

Un árbol es un grafo conexo sin ciclos.

Equivalentemente, es un grafo en el que entre cualquier par de vértices existe un único camino.

Para un árbol se cumple:

\[
|A| = |V| - 1
\]

\subsection*{Grafos planares}

Un grafo es \textbf{planar} si puede dibujarse en el plano sin que sus aristas se crucen, excepto en los vértices.

Un dibujo de un grafo planar sin cruces se llama \textbf{representación plana} o \textbf{incrustación plana}.

Las regiones del plano delimitadas por las aristas se llaman \textbf{caras}.  
Existe siempre una cara exterior no acotada.

\begin{figure}[htbp]
\centering
\begin{tikzpicture}[scale=1]
\node[v] (A) at (0,0) {};
\node[v] (B) at (2,0) {};
\node[v] (C) at (2,2) {};
\node[v] (D) at (0,2) {};
\draw (A)--(B)--(C)--(D)--(A);
\draw (A)--(C);
\draw (B)--(D);
\node at (1,-0.6) {\small Representación planar de $K_4$};
\end{tikzpicture}
\caption{Ejemplo de grafo planar (se puede dibujar sin cruces).}
\end{figure}

\subsection*{Característica de Euler}

Para todo grafo planar conexo se cumple:

\[
|V| - |A| + |F| = 2
\]

donde:

\begin{itemize}
    \item $|V|$ = número de vértices
    \item $|A|$ = número de aristas
    \item $|F|$ = número de caras
\end{itemize}

Esta relación se conoce como \textbf{característica de Euler}.

\begin{figure}[htbp]
\centering
\begin{tikzpicture}[scale=1]
\node[v] (A) at (0,0) {};
\node[v] (B) at (3,0) {};
\node[v] (C) at (3,2) {};
\node[v] (D) at (0,2) {};
\node[v] (E) at (1.5,1) {};

\draw (A)--(B)--(C)--(D)--(A);
\draw (A)--(E)--(B);
\draw (B)--(E)--(C);
\draw (C)--(E)--(D);
\draw (D)--(E)--(A);

% Etiquetas opcionales de caras (solo ilustrativo)
\node at (1.5,1.65) {\tiny $F_1$};
\node at (2.35,1.0) {\tiny $F_2$};
\node at (1.5,0.35) {\tiny $F_3$};
\node at (0.65,1.0) {\tiny $F_4$};
\node at (1.5,-0.35) {\tiny exterior};

\end{tikzpicture}
\caption{Ejemplo planar con varias caras (incluida la exterior) para ilustrar $|V|-|A|+|F|=2$.}
\end{figure}

\subsection*{Consecuencias de la característica de Euler}

Para grafos planares simples con $|V| \ge 3$ se cumple:

\[
|A| \le 3|V| - 6
\]

Si además el grafo es bipartito (no tiene ciclos de longitud impar):

\[
|A| \le 2|V| - 4
\]

Estas desigualdades permiten demostrar que ciertos grafos no son planares.

\subsection*{Grafos no planares fundamentales}

Dos grafos son especialmente importantes:

\begin{itemize}
    \item El grafo completo $K_5$.
    \item El grafo bipartito completo $K_{3,3}$.
\end{itemize}

Ninguno de ellos es planar.

El teorema de Kuratowski establece que un grafo es no planar si y solo si contiene una subdivisión de $K_5$ o $K_{3,3}$.

\subsection*{Subdivisión de una arista}

Una subdivisión consiste en sustituir una arista por un camino formado por nuevos vértices de grado 2.

Dos grafos se consideran equivalentes respecto a planaridad si uno puede obtenerse del otro mediante subdivisiones.

\subsection*{Dual de un grafo planar}

Dado un grafo planar conexo, se puede construir su \textbf{grafo dual}:

\begin{itemize}
    \item Se coloca un vértice en cada cara del grafo original.
    \item Se unen dos vértices del dual si sus caras correspondientes comparten una arista.
\end{itemize}

Cada arista del grafo original corresponde a una arista del dual.

\begin{figure}[htbp]
\centering
\begin{tikzpicture}[scale=1]
% Grafo primal: cuadrado con diagonal
\node[v,label=below left:$a$] (a) at (0,0) {};
\node[v,label=below right:$b$] (b) at (3,0) {};
\node[v,label=above right:$c$] (c) at (3,2) {};
\node[v,label=above left:$d$] (d) at (0,2) {};
\draw (a)--(b)--(c)--(d)--(a);
\draw (a)--(c);

% "Centros" de caras (dual)
\node[v,fill=white] (f1) at (1.1,0.7) {}; % cara triángulo abc
\node[v,fill=white] (f2) at (1.9,1.3) {}; % cara triángulo acd
\node[v,fill=white] (fext) at (1.5,2.8) {}; % exterior (simplificado)

% Aristas del dual (en gris)
\begin{scope}[every edge/.style={draw, thick, dashed}]
\draw (f1)--(f2);     % cruzando diagonal a-c
\draw (f1)--(fext);   % una conexión hacia exterior
\draw (f2)--(fext);   % otra hacia exterior
\end{scope}

\node at (1.5,-0.6) {\small Primal (negro) y dual (trazado)};
\end{tikzpicture}
\caption{Idea de dualidad: un vértice por cara y una arista por cada arista compartida entre caras.}
\end{figure}

\subsection*{Coloración de grafos}

Una \textbf{coloración} asigna colores a los vértices de un grafo de forma que vértices adyacentes tengan colores distintos.

El número mínimo de colores necesarios se llama \textbf{número cromático} y se denota $\chi(G)$.

\subsection*{Teorema de los cuatro colores}

Todo grafo planar puede colorearse con a lo sumo cuatro colores.

Equivalentemente, cualquier mapa plano puede colorearse con cuatro colores sin que regiones adyacentes compartan color.

\subsection*{Coloración de aristas}

También se pueden colorear aristas de modo que aristas incidentes en un mismo vértice tengan colores distintos.

El mínimo número de colores necesarios se denomina \textbf{índice cromático}.

\subsection*{Grafos dirigidos}

Un grafo dirigido o \textbf{digrafo} es un grafo cuyas aristas tienen orientación.

Se representan mediante pares ordenados $(u,v)$.

En un digrafo se distinguen:

\begin{itemize}
    \item \textbf{Grado de entrada} de un vértice: número de aristas que llegan a él.
    \item \textbf{Grado de salida}: número de aristas que parten de él.
\end{itemize}

\begin{figure}[htbp]
\centering
\begin{tikzpicture}[>=Latex, scale=1]

\node[v,label=left:$u$] (u) at (0,1) {};
\node[v,label=right:$v$] (v) at (2,1) {};
\node[v,label=below:$w$] (w) at (1,-0.5) {};

\draw[->] (u)--(v);
\draw[->] (u)--(w);
\draw[->] (w)--(v);

% Lazo dirigido en v (corregido)
\draw[->] (v) to[out=45,in=-45,looseness=8] (v);

\node at (1,-1.2) {\small Ejemplo de digrafo};

\end{tikzpicture}
\caption{En digrafos: grado de salida (aristas que salen) y de entrada (aristas que entran).}
\end{figure}

\subsection*{Caminos en digrafos}

Un camino dirigido es una sucesión de vértices tal que cada arista se recorre en el sentido indicado.

Un digrafo es \textbf{fuertemente conexo} si existe un camino dirigido entre cualquier par ordenado de vértices.

\subsection*{Matrices asociadas a grafos}

\textbf{Matriz de adyacencia:}

Es una matriz $A = (a_{ij})$ donde

\[
a_{ij} =
\begin{cases}
1 & \text{si existe arista entre } v_i \text{ y } v_j, \\
0 & \text{en caso contrario.}
\end{cases}
\]

En grafos no dirigidos simples es simétrica.

\textbf{Matriz de incidencia:}

Es una matriz que relaciona vértices y aristas.  
Cada columna representa una arista y cada fila un vértice.

\subsection*{Componentes conexas}

Una componente conexa es un subgrafo maximal conexo.

Todo grafo puede descomponerse como unión disjunta de sus componentes conexas.

\subsection*{Subgrafos}

Un subgrafo de $G=(V,A)$ es un grafo $H=(V',A')$ tal que:

\[
V' \subseteq V, \qquad A' \subseteq A.
\]

Si contiene todas las aristas entre los vértices de $V'$ presentes en $G$, se denomina \textbf{subgrafo inducido}.

\section*{Algoritmos de grafos}

\subsection*{Árbol recubridor}

Un \textbf{árbol recubridor} de un grafo $G$ es un subgrafo que:

\begin{itemize}
    \item Contiene todos los vértices de $G$
    \item Es conexo
    \item No tiene ciclos (es un árbol)
\end{itemize}

Un mismo grafo puede tener múltiples árboles recubridores.

\subsection*{Búsqueda en profundidad (DFS)}

La \textbf{Depth First Search (DFS)} recorre el grafo explorando lo más profundo posible antes de retroceder.

\subsubsection*{Idea del algoritmo}

\begin{enumerate}
    \item Elegir un vértice inicial.
    \item Marcarlo como visitado.
    \item Ir a un vértice adyacente no visitado.
    \item Repetir hasta no poder avanzar.
    \item Retroceder al último vértice con vecinos no visitados.
\end{enumerate}

\textbf{Propiedades:}

\begin{itemize}
    \item Genera un árbol recubridor.
    \item Depende del vértice inicial y del orden de los vecinos.
    \item Útil para detectar componentes conexas, ciclos y puentes.
    \item Complejidad: $O(|V| + |A|)$
\end{itemize}

\subsection*{Búsqueda en anchura (BFS)}

La \textbf{Breadth First Search (BFS)} recorre el grafo por niveles.

\subsubsection*{Idea del algoritmo}

\begin{enumerate}
    \item Elegir un vértice inicial.
    \item Marcarlo como visitado.
    \item Visitar todos sus vecinos.
    \item Después los vecinos de esos vecinos.
\end{enumerate}

Se implementa con una cola.

\textbf{Propiedades:}

\begin{itemize}
    \item Genera un árbol recubridor por niveles.
    \item Permite calcular distancias mínimas en grafos no ponderados.
    \item Complejidad: $O(|V| + |A|)$
\end{itemize}

\subsection*{DFS vs BFS}

\begin{itemize}
    \item DFS: busca caminos profundos (``autopista larga'').
    \item BFS: explora uniformemente por niveles.
\end{itemize}

Aplicaciones principales:

\begin{itemize}
    \item DFS: componentes conexas, ciclos, puentes.
    \item BFS: caminos mínimos en grafos no ponderados.
\end{itemize}

\subsection*{Componentes conexas}

Un grafo es conexo si existe un camino entre cualquier par de vértices.

Para hallar las componentes conexas:

\begin{enumerate}
    \item Ejecutar DFS o BFS desde un vértice no visitado.
    \item Marcar todos los vértices alcanzados.
    \item Repetir con otro vértice no visitado.
\end{enumerate}

Cada ejecución produce una componente conexa.

\subsection*{Puentes y vértices de corte}

\textbf{Puente:} arista cuya eliminación aumenta el número de componentes conexas.

\textbf{Vértice de corte:} vértice cuya eliminación aumenta el número de componentes conexas.

Son elementos críticos para la conectividad del grafo.

\subsection*{Distancias en grafos}

\begin{itemize}
    \item Distancia entre dos vértices: longitud del camino más corto entre ellos.
    \item Excentricidad de un vértice: mayor distancia a cualquier otro vértice.
    \item Radio del grafo: mínimo de las excentricidades.
    \item Diámetro del grafo: máximo de las excentricidades.
\end{itemize}

Se calculan mediante BFS en grafos no ponderados.

\subsection*{Árbol recubridor mínimo}

En un grafo ponderado, es un árbol recubridor con peso total mínimo.

\subsubsection*{Algoritmo de Kruskal}

\begin{enumerate}
    \item Ordenar las aristas por peso creciente.
    \item Añadir la arista más ligera que no forme ciclo.
    \item Repetir hasta obtener $|V|-1$ aristas.
\end{enumerate}

\subsubsection*{Algoritmo de Prim}

\begin{enumerate}
    \item Elegir un vértice inicial.
    \item Añadir la arista de menor peso que conecte el árbol con un vértice externo.
    \item Repetir hasta cubrir todos los vértices.
\end{enumerate}

\subsection*{Caminos mínimos — Algoritmo de Dijkstra}

Permite encontrar la distancia mínima desde un vértice origen a todos los demás en grafos con pesos no negativos.

\subsubsection*{Idea del algoritmo}

\begin{enumerate}
    \item Inicializar todas las distancias a $\infty$, excepto el origen ($0$).
    \item Seleccionar el vértice no visitado con menor distancia.
    \item Actualizar las distancias a sus vecinos.
    \item Marcar el vértice como visitado.
    \item Repetir hasta visitar todos.
\end{enumerate}

\textbf{Complejidad:} $O(|V|^2)$ (implementación básica).

\subsection*{Flujo máximo}

Se modela el grafo como una red donde cada arista tiene capacidad.

Objetivo: hallar el flujo máximo entre un origen $s$ y un destino $t$.

\subsection*{Corte mínimo}

Un \textbf{corte} es un conjunto de aristas cuya eliminación desconecta $s$ de $t$.

El peso del corte es la suma de capacidades de dichas aristas.

\textbf{Teorema max-flow min-cut:}

\[
\text{Flujo máximo} = \text{peso del corte mínimo}
\]

\subsection*{Algoritmo de Ford–Fulkerson}

Método para calcular el flujo máximo.

\subsubsection*{Idea}

\begin{enumerate}
    \item Inicializar flujo cero.
    \item Buscar un camino aumentante de $s$ a $t$.
    \item Aumentar el flujo según la capacidad mínima del camino.
    \item Actualizar el grafo residual.
    \item Repetir hasta que no existan más caminos.
\end{enumerate}
\end{document}